\section{A source-closed finite-prefix certificate}
\label{app:source-closed-prefix}

This appendix proves the finite-prefix input used by
Section~\ref{sec:affine-provider}.  The theorem-level datum is only a cofinal
family of homogeneous affine covers with a uniform positive value margin and
uniformly thin incidence against all late Tail windows.  The lcm coordinate,
the fixed root identities and the first few bridge rows are certificate data.

\subsection{A fixed finite root fibre}

Put
\[
P(m)=\frac1m+\frac1{m+1}.
\]
The empty state together with the following five strict binary states gives a
complete affine fibre of the subgroup of order six:
\[
\frac16=
\sum_{m\in\{35,56,90,119,152,170,455,494,714\}}P(m),
\]
\[
\frac13=
\sum_{m\in\{17,35,44,84,104,119,132,170,272,285,455,494,560,714,935\}}P(m),
\]
\[
\frac12=
\sum_{m\in\{9,17,34,44,77,84,90\}}P(m),
\]
\[
\frac23=
\sum_{m\in\{7,16,34,44,51,65,77,119,132,153,170,272,285,455,494\}}P(m),
\]
\[
\frac56=P(2).
\tag{SC1}\label{eq:fixed-root-identities}
\]
Their grades are respectively
\[
0,9,15,7,15,1.
\tag{SC2}\label{eq:fixed-root-grades}
\]
Each identity in \eqref{eq:fixed-root-identities} is an exact rational
identity, and every displayed start set has gaps at least three.  Thus the six
states are legal and their residue image is the full affine $H_6$-fibre.  The
largest root mass is $5/6$.

The only bridge starts requiring a finite compatibility check before the
uniform scale estimate applies are
\[
12;\qquad 20,30,60;\qquad 140,210,420;\qquad 840.
\tag{SC3}\label{eq:early-bridge-starts}
\]
Direct comparison with the finite root support shows that every one of these
bridge envelopes is disjoint and non-adjacent to every root envelope.  The
next nontrivial bridge has smallest start $1260$, while the root support ends
at $936$; from that point onward compatibility follows from the general scale
estimate below.  No exact sum of all finite masses is used.

\subsection{Scale-proper simple bridges}

Let
\[
L_n=\operatorname{lcm}(1,\ldots,n),
\qquad
E_n=H_{L_n}.
\]
At a nontrivial prime-power jump $n=p^e$ put
\[
D=L_{n-1},
\qquad L_n=pD.
\]
The quotient $E_n/E_{n-1}$ has order $p$.

For the finitely many jumps through $n=8$ use the coefficient rows encoded in
\eqref{eq:early-bridge-starts}.  For every later jump let
\[
R_p=\{1,2,4,\ldots,2^{m-1}\},
\qquad m=\lceil\log_2p\rceil.
\]
Every coefficient is $<p$, and the subset sums of $R_p$ contain every integer
from $0$ through $p-1$.  Hence reduction modulo $p$ is surjective.  Moreover
\[
|R_p|=O(\log p),
\qquad
\sum_{r\in R_p}r=2^m-1<2p.
\tag{SC4}\label{eq:binary-coefficient-row}
\]
Since every $r\in R_p$ satisfies $r<p\le n$, it divides $D=L_{n-1}$.

For $r\in R_p$ set
\[
a_r=\frac{pD}{r}
\]
and use the endpoint switch
\[
[a_r+1,a_r+2]
\longleftrightarrow
[a_r,a_r+2].
\tag{SC5}\label{eq:source-closed-bridge-switch}
\]
The alternatives both have grade one and their reciprocal-value difference is
$r/(pD)$.  Thus the Boolean row is a homogeneous relative affine cover of the
new simple quotient.

If $r<s$ are powers of two, then $s\ge2r$, and therefore
\[
|a_r-a_s|
=\frac{pD(s-r)}{rs}
\ge\frac{pD}{s}
>D.
\]
Thus every later bridge row is internally strict.

Successive rows are separated just as directly.  Let $D'=L_n$ be the upper
modulus of one bridge and hence the lower modulus of the next nontrivial jump
$n'=q^f$.  The first row lies below $D'+2$.  If $s<q$ is a coefficient in the
next row, then
\[
\frac{qD'}s-D'
=\frac{D'(q-s)}s
\ge\frac{D'}{q-1}.
\]
If $f=1$, then $q\ge11$ outside the finite regime and the coprime integers
$q-1,q-2$ both divide $D'$, so $D'/(q-1)\ge q-2\ge9$.  If $f\ge2$, then
$q^f-1\mid D'$ and
\[
\frac{D'}{q-1}
\ge\frac{q^f-1}{q-1}
=1+q+\cdots+q^{f-1}\ge4
\]
for $q^f\ge9$.  Hence the next complete envelope starts at least two integers
after the preceding complete envelope ends.  Together with the finite check
\eqref{eq:early-bridge-starts}, all bridge rows are mutually compatible with
one another and with the fixed root.

Each state of one late row has mass
\[
W(\text{row})
<\frac3{pD}\sum_{r\in R_p}r
<\frac6D.
\tag{SC6}\label{eq:source-closed-row-load}
\]
The nontrivial lower moduli grow by at least a factor two.  From the first
fully uniform lower modulus $420$ onward, the entire future bridge load is
therefore less than
\[
\frac{12}{420}=\frac1{35}.
\tag{SC7}\label{eq:source-closed-tail-load}
\]
For the three earlier nontrivial rows one may use the direct bounds
\[
\frac14,\qquad \frac3{10},\qquad \frac1{20}.
\]
Combining these with the root bound $5/6$ and
\eqref{eq:source-closed-tail-load} leaves a fixed positive distance from $2$.

\subsection{Homogenization with fixed complexity}

Let
\[
M=\max\{0,9,15,7,15,1\}.
\]
For each truncation $B$, first form the physical product of the fixed root with
all bridge rows through $B$.  The bridge rows are homogeneous, so the six raw
branch grades differ only by the fixed root grades.

Use $M$ remote copies of the exact grade-neutral identity
\[
P(a)+P(a^2+3a+1)
=
P(a+2)+P\!\left(\frac{a(a+3)}2\right)+P(a(a+3)).
\tag{SC8}\label{eq:source-closed-grade-neutral}
\]
For sufficiently large $a$ the two sides are legal, have equal reciprocal
value and grades two and three.  Construct the $M$ coordinates successively
beyond the complete raw prefix, making their total common mass less than
$1/10$.  For a root branch of grade $g$, take the high alternative on exactly
$M-g$ coordinates.  Every final branch then has one common grade, while all
branches receive the same exact-value translation.

Using the coarse bounds above,
\[
W
<
\frac56+\frac14+\frac3{10}+\frac1{20}+\frac1{35}+\frac1{10}
<\frac53.
\tag{SC9}\label{eq:source-closed-slack}
\]
Hence the entire moving family has a uniform positive literalization margin;
for example $W<2-1/3$.

The number and component complexity of the grade-neutral coordinates are
fixed independently of $B$.  Their numerical positions may move with $B$, but
their contribution to every Tail-window envelope count is uniformly $O(1)$.

\subsection{Uniform source closure}

Let $I_Q(a,A)$ be a fixed enlargement of a Tail current window,
\[
I_Q(a,A)=
\left[a\frac{Q^2}{\log Q},AQ^2\right]
\qquad(0<a<A).
\]
Suppose a late bridge row at jump $n=p^e$, lower modulus $D=L_{n-1}$, meets
$I_Q(a,A)$.  Since its bases lie between $D$ and $pD$, the upper edge gives
\[
D=O_{a,A}(Q^2).
\]
The primorial below $n$ divides $D$, so the Chebyshev lower bound from
Appendix~\ref{app:prime-supply} gives
\[
n=O_{a,A}(\log Q),
\qquad p=O_{a,A}(\log Q).
\]
The lower edge of the Tail window then gives
\[
D\gg_{a,A}\frac{Q^2}{(\log Q)^2}.
\]
Only $O_{a,A}(\log\log Q)$ nontrivial lower moduli can lie in this
multiplicative window, and each bridge row has $O(\log p)=O(\log\log Q)$
envelopes.  The fixed root and fixed-number neutral equalizers contribute only
$O(1)$.  Consequently, uniformly over every truncation $B$,
\[
\sup_B N_Q(\operatorname{Prefix}_B)
=O_{a,A}((\log\log(3Q))^2+1)
=o(Q/\log Q).
\tag{SC10}\label{eq:uniform-source-closure}
\]

This is the source-closed estimate.  Restrict the numerical endpoint indices
to the proper cofinal object of eligible Tail bases.  For every sufficiently
late eligible $B$, the same family covers a homogeneous affine correction
\[
0\longrightarrow E_B
\]
with the uniform margin \eqref{eq:source-closed-slack} and is a uniformly
row-transparent obstacle at that very same source $E_B$, including the
boundary-predecessor rows.  Thus the Prefix source map may be taken to be the
identity on this cofinal eligible-base object, and is automatically
$\operatorname{Tendsto}$ to the eligible-source filter.

The explicit root identities have now reached their last consumer.  Above this
appendix the Prefix is used only through its homogeneous affine cover, uniform
positive margin and source-closed transparency.